\section{Explicit two-point minimax lower bounds}
\label{sec:minimax-lower-bound}

Let \(P_\theta^{(T)}\) denote the law of the observed trajectory in the
capped-geometric Gaussian submodel of
Section~\ref{sec:explicit-capped-gaussian}.  The data length remains \(T\)
and the target horizon remains \(N\).

For fixed \(p^\circ=(p_1^\circ,p_2^\circ)\in(p_-,p_+)^2\), let
\(Q^\circ\) and \(\eta^\circ\) be the transition matrix and stationary law
of the augmented age chain.  Put
\[
 y(i,r)=2\mathbf 1_{\{i=2\}}-1,\qquad
 \alpha^\circ=(\eta^\circ)^\top y,\qquad
 \bar y=y-\alpha^\circ\mathbf 1,
\]
\[
 \Pi^\circ=\mathbf 1(\eta^\circ)^\top,\qquad
 R^\circ=Q^\circ-\Pi^\circ,
\]
and use the stationary inner product
\(\langle u,v\rangle_{\eta^\circ}
=u^\top\operatorname{diag}(\eta^\circ)v\).  Define
\begin{align}
 A_N^\circ
 &:=
 \operatorname{Var}_{p^\circ}\!\left(\sum_{t=1}^N Y_t\right)
 \nonumber\\
 &=N\langle\bar y,\bar y\rangle_{\eta^\circ}
 +2\left\langle
 \bar y,\,
 R^\circ\{(N-1)I-NR^\circ+(R^\circ)^N\}
 (I-R^\circ)^{-2}\bar y
 \right\rangle_{\eta^\circ}.
 \label{eq:lb-AN}
\end{align}
If \(\tau_i^\circ=\tau_D(p_i^\circ)\) and
\(\nu_i^{\circ\,2}=\operatorname{Var}\{\min(G_i,D)\}\), set
\begin{equation}
 A_\infty^\circ
 =
 4\,
 \frac{(\tau_2^\circ)^2\nu_1^{\circ\,2}
       +(\tau_1^\circ)^2\nu_2^{\circ\,2}}
      {(\tau_1^\circ+\tau_2^\circ)^3}.
 \label{eq:lb-Ainfty}
\end{equation}
The two target functionals along the amplitude submodel are therefore
\begin{equation}
 G_N(a,p^\circ)=N\sigma^2+a^2A_N^\circ,\qquad
 v_\infty(a,p^\circ)=\sigma^2+a^2A_\infty^\circ.
 \label{eq:lb-target-amplitude}
\end{equation}

\subsection{An explicit relative-entropy envelope}

For \(p\in(0,1)\), define
\[
 \tau_D(p)=\sum_{r=0}^{D-1}(1-p)^r,\qquad
 e_D(p)=\sum_{r=0}^{D-2}(1-p)^r,
\]
\[
 c_D(p)=\sum_{r=0}^{D-1}r(1-p)^r
 =\frac{(1-p)-D(1-p)^D+(D-1)(1-p)^{D+1}}{p^2},
\]
and
\[
 \operatorname{kl}(u\|v)
 =u\log\frac uv+(1-u)\log\frac{1-u}{1-v}.
\]

\begin{lemma}[Complete-data entropy envelope]
\label{lem:lb-complete-kl}
For \(\theta=(a,p_1,p_2)\) and
\(\theta'=(a',p_1',p_2')\), write
\(T_c(\theta)=\tau_D(p_1)+\tau_D(p_2)\).  Then
\begin{align}
 D_{\mathrm{KL}}\!\left(P_\theta^{(T)}\middle\|P_{\theta'}^{(T)}\right)
 \leq \mathcal K_T(\theta,\theta')
 &:=
 \log\frac{T_c(\theta')}{T_c(\theta)}
 +\frac{1}{T_c(\theta)}
 \sum_{i=1}^2c_D(p_i)
 \log\frac{1-p_i}{1-p_i'}
 \nonumber\\
 &\quad
 +(T-1)\sum_{i=1}^2
 \frac{e_D(p_i)}{T_c(\theta)}
 \operatorname{kl}(p_i\|p_i')
 +\frac{T(a-a')^2}{2\sigma^2}.
 \label{eq:lb-complete-kl}
\end{align}
\end{lemma}

\begin{proof}
The stationary law of the augmented age chain is
\[
 \eta_\theta(i,r)=
 \frac{(1-p_i)^{r-1}}{T_c(\theta)},\qquad 1\leq r\leq D.
\]
Consequently,
\[
 D_{\mathrm{KL}}(\eta_\theta\|\eta_{\theta'})
 =
 \log\frac{T_c(\theta')}{T_c(\theta)}
 +\frac{1}{T_c(\theta)}
 \sum_{i=1}^2c_D(p_i)
 \log\frac{1-p_i}{1-p_i'}.
\]
At an age \(r<D\), the transition row is Bernoulli with switching
probability \(p_i\); at \(r=D\), the exit is forced under both parameters.
Stationarity thus makes the transition contribution equal to
\[
 (T-1)\sum_{i=1}^2
 \frac{e_D(p_i)}{T_c(\theta)}
 \operatorname{kl}(p_i\|p_i').
\]
The conditional Gaussian contribution is
\[
 \sum_{t=1}^T
 D_{\mathrm{KL}}\{N(aY_t,\sigma^2)\|N(a'Y_t,\sigma^2)\}
 =\frac{T(a-a')^2}{2\sigma^2}.
\]
The chain rule gives equality for the complete trajectory
\((Z_{1:T},X_{1:T})\); projection onto \(X_{1:T}\) proves
\eqref{eq:lb-complete-kl}.
\end{proof}

\subsection{Finite-sample two-point bound}

\begin{theorem}[Explicit minimax lower bound]
\label{thm:lb-two-point}
Fix \(a_0\in[a_-,a_+)\), \(p^\circ\in(p_-,p_+)^2\), and
\(0<d\leq a_+-a_0\).  Let
\(\theta_0=(a_0,p^\circ)\) and
\(\theta_1=(a_0+d,p^\circ)\).  Every estimator based on \(X_{1:T}\)
satisfies
\begin{align}
 \max_{j=0,1}
 P_{\theta_j}^{(T)}
 \left\{\|\widehat\theta_T-\theta_j\|_2\geq\frac d2\right\}
 &\geq
 \frac14\exp\!\left\{-\frac{Td^2}{2\sigma^2}\right\},
 \label{eq:lb-theta}\\
 \max_{j=0,1}
 P_{\theta_j}^{(T)}
 \left\{\left|\widehat G_{N,T}-G_N(\theta_j)\right|
 \geq\frac{A_N^\circ}{2}(2a_0d+d^2)\right\}
 &\geq
 \frac14\exp\!\left\{-\frac{Td^2}{2\sigma^2}\right\},
 \label{eq:lb-GN}\\
 \max_{j=0,1}
 P_{\theta_j}^{(T)}
 \left\{\left|\widehat v_{\infty,T}-v_\infty(\theta_j)\right|
 \geq\frac{A_\infty^\circ}{2}(2a_0d+d^2)\right\}
 &\geq
 \frac14\exp\!\left\{-\frac{Td^2}{2\sigma^2}\right\}.
 \label{eq:lb-vinfty}
\end{align}
\end{theorem}

\begin{proof}
Lemma~\ref{lem:lb-complete-kl}, with \(p^\circ\) fixed, gives
\[
 D_{\mathrm{KL}}\!\left(
 P_{\theta_0}^{(T)}\middle\|P_{\theta_1}^{(T)}
 \right)
 \leq \frac{Td^2}{2\sigma^2}.
\]
The nearest-point test associated with any estimator and the
Bretagnolle--Huber inequality \cite{BretagnolleHuber1978} give
\[
 \max_{j=0,1}P_{\theta_j}^{(T)}
 \left\{d_g(\widehat g_T,g(\theta_j))
 \geq\frac{d_g(g(\theta_0),g(\theta_1))}{2}\right\}
 \geq
 \frac14
 \exp\!\left\{
 -D_{\mathrm{KL}}(P_{\theta_0}^{(T)}\|P_{\theta_1}^{(T)})
 \right\}.
\]
Here \(d_g\) is Euclidean distance for \(g(\theta)=\theta\) and absolute
distance for the scalar functionals.
Equation~\eqref{eq:lb-target-amplitude} supplies the three separations in
\eqref{eq:lb-theta}--\eqref{eq:lb-vinfty}.
\end{proof}

\begin{corollary}[Necessary confidence radii and matching rate]
\label{cor:lb-confidence}
For \(0<\delta<1/8\), put
\[
 d_{T,\delta}
 =\sigma\sqrt{\frac{2}{T}\log\frac{1}{8\delta}},
 \qquad
 y_\delta=\log\frac{10}{\delta},
 \qquad M=T-2,
\]
and assume \(d_{T,\delta}\leq a_+-a_0\).  If
\(\mathfrak r_{T,\delta}(g)\) denotes the smallest uniform
\((1-\delta)\)-confidence radius for \(g\), namely
\[
 \mathfrak r_{T,\delta}(g)
 =
 \inf_{\widehat g}\inf\left\{r>0:
 \sup_{\theta\in\Theta_{\rm cg}}
 P_\theta^{(T)}
 \{\ell_g(\widehat g,g(\theta))>r\}\leq\delta\right\},
\]
where \(\ell_\theta\) is Euclidean distance and the two functional losses
are absolute distance, then
\begin{align}
 \mathfrak r_{T,\delta}(\theta)
 &\geq\frac{d_{T,\delta}}2,
 \label{eq:lb-radius-theta}\\
 \mathfrak r_{T,\delta}(G_N)
 &\geq A_N^\circ
 \left(a_0d_{T,\delta}+\frac{d_{T,\delta}^2}{2}\right),
 \label{eq:lb-radius-GN}\\
 \mathfrak r_{T,\delta}(v_\infty)
 &\geq A_\infty^\circ
 \left(a_0d_{T,\delta}+\frac{d_{T,\delta}^2}{2}\right).
 \label{eq:lb-radius-vinfty}
\end{align}
For the estimator of Section~\ref{sec:explicit-capped-gaussian}, define
\[
 U_{T,\delta}
 =
 2K_{\rm id,cg}
 \left[
 \frac{40y_\delta}{M\underline g_3}
 +\sqrt{\frac{8y_\delta}{M\underline g_3}
 \left(1+\frac{1}{M\underline g_3}\right)}
 \right].
\]
Then
\begin{equation}
 \mathfrak r_{T,\delta}(\theta)\leq U_{T,\delta},
 \qquad
 \mathfrak r_{T,\delta}(g)\leq L_gU_{T,\delta},
 \quad g\in\{G_N,v_\infty\},
 \label{eq:lb-upper-radius}
\end{equation}
where \(L_g\) is the Lipschitz constant of \(g\) on
\(\Theta_{\rm cg}\).  Hence, for a fixed interior parameter box,
\[
 \mathfrak r_{T,\delta}(\theta)
 \asymp
 \sqrt{\frac{\log(1/\delta)}{T}},
 \qquad
 \mathfrak r_{T,\delta}(v_\infty)
 \asymp
 \sqrt{\frac{\log(1/\delta)}{T}},
\]
up to structural constants.  Moreover,
\begin{equation}
 \frac{A_N^\circ}{N}\longrightarrow A_\infty^\circ>0,
 \label{eq:lb-AN-limit}
\end{equation}
so the rate for the unnormalised target is
\[
 \mathfrak r_{T,\delta}(G_N)
 \asymp
 N\sqrt{\frac{\log(1/\delta)}{T}}
\]
whenever \(N\) is in the regime where the uniform \(O(N)\) Lipschitz
bound applies.
\end{corollary}

\begin{proof}
With \(d=d_{T,\delta}\), the right side of
\eqref{eq:lb-theta}--\eqref{eq:lb-vinfty} equals \(2\delta\), proving
\eqref{eq:lb-radius-theta}--\eqref{eq:lb-radius-vinfty}.  Setting
\(x=\log(2/\delta)\) in the upper bound of
Section~\ref{sec:explicit-capped-gaussian}, and using
\(\sqrt{u+v}\leq\sqrt u+\sqrt v\), gives
\eqref{eq:lb-upper-radius}.  Finally,
\[
 \frac{A_N^\circ}{N}
 =\operatorname{Var}(Y_0)
 +2\sum_{h=1}^{N-1}
 \left(1-\frac hN\right)\operatorname{Cov}(Y_0,Y_h)
 \longrightarrow A_\infty^\circ
\]
by the geometric covariance bound.  On the compact parameter box, the
covariances and their first parameter derivatives are uniformly summable,
which gives the corresponding \(O(N)\) upper Lipschitz bound.
\end{proof}

\subsection{Persistence-sensitive lower bound}

On the symmetric submodel \(p_1=p_2=p\), define
\[
 \omega_D(p)=\frac{e_D(p)}{\tau_D(p)},\qquad
 \nu_D^2(p)=
 \sum_{r=0}^{D-1}(2r+1)(1-p)^r-\tau_D(p)^2,
\]
\[
 H_D(p)=\frac{\nu_D^2(p)}{\tau_D(p)},\qquad
 v_\infty(a,p,p)=\sigma^2+a^2H_D(p),
\]
and
\[
 K_D^{\rm age}(p,q)
 =
 \log\frac{\tau_D(q)}{\tau_D(p)}
 +\frac{c_D(p)}{\tau_D(p)}
 \log\frac{1-p}{1-q}.
\]
All derivatives in the local bound below are explicit finite sums:
\[
 \tau_D'(p)=-\sum_{r=1}^{D-1}r(1-p)^{r-1},
\]
\[
 \{\nu_D^2\}'(p)
 =-\sum_{r=1}^{D-1}r(2r+1)(1-p)^{r-1}
 -2\tau_D(p)\tau_D'(p),
\]
\[
 H_D'(p)
 =\frac{\{\nu_D^2\}'(p)\tau_D(p)
 -\nu_D^2(p)\tau_D'(p)}{\tau_D(p)^2}.
\]

\begin{proposition}[Explicit persistence lower bound]
\label{prop:lb-persistence}
For interior \(p,q\in[p_-,p_+]\), every estimator satisfies
\begin{align}
 &\max_{\vartheta\in\{p,q\}}
 P_{(a,\vartheta,\vartheta)}^{(T)}
 \left\{
 \left|\widehat v_{\infty,T}
 -v_\infty(a,\vartheta,\vartheta)\right|
 \geq
 \frac{a^2}{2}|H_D(q)-H_D(p)|
 \right\}
 \nonumber\\
 &\hspace{35mm}\geq
 \frac14\exp\!\left[
 -K_D^{\rm age}(p,q)
 -(T-1)\omega_D(p)\operatorname{kl}(p\|q)
 \right].
 \label{eq:lb-persistence-finite}
\end{align}
If \(q_T=p+h/\sqrt T\) and \(H_D'(p)\neq0\), then
\begin{align}
 K_D^{\rm age}(p,q_T)
 +(T-1)\omega_D(p)\operatorname{kl}(p\|q_T)
 &\longrightarrow
 \frac{h^2\omega_D(p)}{2p(1-p)},
 \label{eq:lb-persistence-kl-limit}\\
 \frac{\sqrt T}{2}
 |v_\infty(a,q_T,q_T)-v_\infty(a,p,p)|
 &\longrightarrow
 \frac{a^2|H_D'(p)h|}{2}.
 \label{eq:lb-persistence-target-limit}
\end{align}
Thus the local unavoidable scale is
\begin{equation}
 a^2|H_D'(p)|
 \sqrt{\frac{p(1-p)}{T\omega_D(p)}}.
 \label{eq:lb-persistence-scale}
\end{equation}
\end{proposition}

\begin{proof}
Apply Lemma~\ref{lem:lb-complete-kl} to
\((a,p,p)\) and \((a,q,q)\), and then apply the two-point testing inequality
used in Theorem~\ref{thm:lb-two-point}.  Equations
\eqref{eq:lb-persistence-kl-limit} and
\eqref{eq:lb-persistence-target-limit} follow from the second-order
expansion
\[
 \operatorname{kl}(p\|p+u)
 =\frac{u^2}{2p(1-p)}+O(u^3)
\]
and first-order differentiation of \(H_D\).
\end{proof}

\subsection{Role of the exact local identification modulus}

Let \(\mathcal I\) be the inverse map in \eqref{eq:cg-inverse}, let
\(J_\theta=D\mathcal I(\Psi(\theta))\), and let
\[
 \Sigma_\Phi(\theta)
 =
 \sum_{k\in\mathbb Z}
 \operatorname{Cov}_\theta\{
 \phi(X_0,X_1,X_2),\phi(X_k,X_{k+1},X_{k+2})\}.
\]

\begin{proposition}[Plug-in fluctuation]
\label{prop:lb-local-kappa}
At an interior parameter, put
\[
 \lambda_\Phi(\theta)
 :=\lambda_{\min}\{\Sigma_\Phi(\theta)\},
\]
and suppose that \(\lambda_\Phi(\theta)>0\).  The direct inverse estimator
\(\widetilde\theta_T=\mathcal I(\widehat z_T)\), locally extended outside
the range of \(\mathcal I\), satisfies
\[
 \sqrt T(\widetilde\theta_T-\theta)
 \ \Longrightarrow\
 N\{0,J_\theta\Sigma_\Phi(\theta)J_\theta^\top\}.
\]
Writing
\[
 K_{\rm id,cg}^{\rm loc}(\theta)
 =\|J_\theta\|_{\infty\to2}
 =\{\kappa_{\rm id,cg}^{\rm loc}(\theta)\}^{-1},
\]
for every \(c>0\),
\begin{equation}
 \liminf_{T\to\infty}
 P_\theta\!\left\{
 \sqrt T\|\widetilde\theta_T-\theta\|_2
 \geq
 cK_{\rm id,cg}^{\rm loc}(\theta)
 \sqrt{\frac{\lambda_\Phi(\theta)}5}
 \right\}
 \geq2\Phi_{\rm N}(-c).
 \label{eq:lb-local-kappa}
\end{equation}
\end{proposition}

\begin{proof}
Geometric ergodicity and boundedness of \(\phi\) give the multivariate
Markov-chain central limit theorem; the displayed weak limit follows by the
delta method.  Since
\[
 \|J_\theta\|_{\infty\to2}
 \leq\sqrt5\,\|J_\theta\|_{2\to2},
\]
\[
 \lambda_{\max}\{
 J_\theta\Sigma_\Phi(\theta)J_\theta^\top\}
 \geq
 \frac{\lambda_\Phi(\theta)}5
 \{K_{\rm id,cg}^{\rm loc}(\theta)\}^2.
\]
Projection onto a maximizing eigenvector gives
\eqref{eq:lb-local-kappa}.
\end{proof}

\begin{remark}[Scope of rate matching]
\label{rem:lb-scope}
Corollary~\ref{cor:lb-confidence} proves minimax-rate optimality in
\(T\), \(\log(1/\delta)\), and \(N\) on a fixed interior box.
Proposition~\ref{prop:lb-local-kappa} shows how the exact local inverse
modulus enters the fluctuation of the specified moment plug-in estimator.
It does not make the global factor
\(K_{\rm id,cg}/\sqrt{\underline g_3}\) a sharp minimax constant.
Indeed, rescaling the feature vector changes \(K_{\rm id,cg}\) while leaving
the observation experiment unchanged; the invariant local object is
\(J_\theta\Sigma_\Phi(\theta)J_\theta^\top\).  Likewise,
\(\underline g_3\) is a certified uniform Doeblin lower bound, whereas
\eqref{eq:lb-persistence-finite} is governed by the explicit local
information coefficient
\(\omega_D(p)/\{p(1-p)\}\).
\end{remark}

\subsection{Quantitative comparison with the closest minimax antecedents}
\label{subsec:lb-exact-literature-map}

\paragraph{Two-state HMM parameter bounds.}
First consider the uncapped Markov benchmark \(D=\infty\), and write

\[
 Q=\begin{pmatrix}1-p&p\\q&1-q\end{pmatrix},\qquad
 s=p+q,\qquad \lambda=1-s.
\]
For \(Y_t=2\mathbf 1_{\{S_t=2\}}-1\), stationarity gives
\[
 B(p,q):=\operatorname{Var}(Y_t)=\frac{4pq}{s^2},\qquad
 \operatorname{Cov}(Y_t,Y_{t+h})=B(p,q)\lambda^h.
\]
Consequently,
\begin{align}
 v_\infty(p,q,a)
 &=\sigma^2+a^2B(p,q)\frac{1+\lambda}{1-\lambda}
 =\sigma^2+\frac{4a^2pq(2-p-q)}{(p+q)^3},
 \label{eq:agn-map-vinfty}\\
 G_N(p,q,a)
 &=N\sigma^2+a^2B(p,q)F_N(\lambda),
 \label{eq:agn-map-GN}
\end{align}
where
\begin{align}
 F_N(\lambda)
 &=N+2\sum_{h=1}^{N-1}(N-h)\lambda^h
 \nonumber\\
 &=N\frac{1+\lambda}{1-\lambda}
 -\frac{2\lambda(1-\lambda^N)}{(1-\lambda)^2},
 \label{eq:agn-map-FN}\\
 F_N'(\lambda)
 &=\frac{2\{(N-1)-(N+1)\lambda
 +(N+1)\lambda^N-(N-1)\lambda^{N+1}\}}
 {(1-\lambda)^3}.
 \label{eq:agn-map-FNprime}
\end{align}

The comparison with \cite{AbrahamGassiatNaulet2023} can be made literal,
despite their finite-alphabet assumption, by applying their estimator to the
observable coarsening \(U_t=\mathbf 1_{\{X_t>0\}}\).  The two Bernoulli
emission vectors are
\[
 f_1=\{\Phi_{\rm N}(a/\sigma),\Phi_{\rm N}(-a/\sigma)\},\qquad
 f_2=\{\Phi_{\rm N}(-a/\sigma),\Phi_{\rm N}(a/\sigma)\},
\]
so their emission-separation coordinate is
\[
 \phi_3=\|f_1-f_2\|_2
 =\sqrt2\{2\Phi_{\rm N}(a/\sigma)-1\}.
\]
With their remaining coordinates
\[
 \phi_1=\frac{q-p}{p+q},\qquad \phi_2=1-p-q=\lambda,
\]
one has
\[
 a=a(\phi_3)
 =\sigma\Phi_{\rm N}^{-1}
 \left(\frac{1+\phi_3/\sqrt2}{2}\right),\qquad
 a'(\phi_3)=\frac{\sigma}
 {2\sqrt2\,\varphi_{\rm N}\{a(\phi_3)/\sigma\}},
\]
where \(\varphi_{\rm N}\) is the standard normal density.  Hence
\begin{align}
 v_\infty(\phi)
 &=\sigma^2+a(\phi_3)^2(1-\phi_1^2)
 \frac{1+\phi_2}{1-\phi_2},
 \label{eq:agn-map-vphi}\\
 G_N(\phi)
 &=N\sigma^2+a(\phi_3)^2(1-\phi_1^2)F_N(\phi_2).
 \label{eq:agn-map-Gphi}
\end{align}
The required delta map is therefore explicit:
\begin{align}
 \partial_{\phi_1}v_\infty
 &=-2a^2\phi_1\frac{1+\phi_2}{1-\phi_2},
 &
 \partial_{\phi_2}v_\infty
 &=\frac{2a^2(1-\phi_1^2)}{(1-\phi_2)^2},
 \nonumber\\
 \partial_{\phi_3}v_\infty
 &=2aa'(\phi_3)(1-\phi_1^2)
 \frac{1+\phi_2}{1-\phi_2},
 \label{eq:agn-map-grad-v}\\
 \partial_{\phi_1}G_N
 &=-2a^2\phi_1F_N(\phi_2),
 &
 \partial_{\phi_2}G_N
 &=a^2(1-\phi_1^2)F_N'(\phi_2),
 \nonumber\\
 \partial_{\phi_3}G_N
 &=2aa'(\phi_3)(1-\phi_1^2)F_N(\phi_2).
 \label{eq:agn-map-grad-G}
\end{align}

In the notation of \cite[Theorem~1]{AbrahamGassiatNaulet2023}, their
transition and emission radii are, up to constants depending on the alphabet
size and on their fixed spectral-gap margin \(L_{\rm A}\),
\begin{equation}
 r_Q(T,x)=
 \frac{x(\delta_{\rm A}\vee\epsilon_{\rm A}\zeta_{\rm A})}
 {\sqrt{T\delta_{\rm A}^2\epsilon_{\rm A}^4\zeta_{\rm A}^6}},
 \qquad
 r_f(T,x)=
 \frac{x}
 {\sqrt{T\delta_{\rm A}^2\epsilon_{\rm A}^4\zeta_{\rm A}^4}},
 \label{eq:agn-raw-rates}
\end{equation}
with tail probability at most \(e^{-x^2}\).  On a fixed interior box,
\(\delta_{\rm A},\epsilon_{\rm A},\zeta_{\rm A},L_{\rm A}\) are fixed and
\(a'(\phi_3)\) is bounded.  Moreover, for
\(|\phi_2|\leq\rho<1\), equations
\eqref{eq:agn-map-FN}--\eqref{eq:agn-map-FNprime} give
\[
 \sup_{|\lambda|\leq\rho}|F_N(\lambda)|\leq C_\rho N,
 \qquad
 \sup_{|\lambda|\leq\rho}|F_N'(\lambda)|
 \leq\frac{2N}{(1-\rho)^2}.
\]
The Jacobian of
\((p,q,f_1,f_2)\mapsto(\phi_1,\phi_2,\phi_3)\) is bounded on this box, so
\[
 \|\widehat\phi-\phi\|_2
 \leq C_{\Box}\{r_Q(T,x)+r_f(T,x)\}
 \leq C_{\Box}\frac{x}{\sqrt T}.
\]
The mean-value theorem applied to
\eqref{eq:agn-map-grad-v}--\eqref{eq:agn-map-grad-G} thus yields
\begin{equation}
 |\widehat v_\infty-v_\infty|
 \leq C_{\Box}\frac{x}{\sqrt T},\qquad
 |\widehat G_N-G_N|
 \leq C_{\Box}N\frac{x}{\sqrt T}.
 \label{eq:agn-delta-conclusion}
\end{equation}
Taking \(x=\sqrt{\log(1/\alpha)}\) reproduces exactly the powers of
\(T\), \(N\), and \(\log(1/\alpha)\) in
Corollary~\ref{cor:lb-confidence}.

Thus the upper-rate statement on a fixed Markov interior box is a direct
functional consequence of the AGN parameter bounds, after the displayed
Bernoulli coarsening; there is no unbounded-gradient or nonlinear obstruction
there.  A raw-parameter lower bound cannot be propagated in reverse by the
delta method: a least-favourable parameter direction may be tangent to a level
set of the functional.  Indeed,
\(\partial_{\phi_1}v_\infty=0\) at \(p=q\).  Nevertheless, our amplitude
pair has
\[
 \partial_av_\infty=2aA_\infty>0,
 \qquad \partial_aG_N=2aA_N>0,
\]
and applying the same two-point functional principle on that regular direction
gives \eqref{eq:lb-GN}--\eqref{eq:lb-vinfty}.  AGN's published lower
theorem is not a formal lower bound for the Gaussian subfamily, because its
least-favourable multinomial pairs need not lie on the Gaussian curve; our
Gaussian KL calculation supplies that missing restriction.  This is a routine
model-specific Le Cam application, not a new minimax rate.  Accordingly, the
fixed-interior \(T^{-1/2}\) and \(NT^{-1/2}\) claims are not presented here as
original consequences.  AGN state their minimax lower radius at fixed
probability \(1/4\); the explicit \(\sqrt{\log(1/\alpha)}\) lower-confidence
calibration in Corollary~\ref{cor:lb-confidence} comes from the
Bretagnolle--Huber calculation above, not from a reverse delta argument.

\begin{remark}[The AGN comparison on the small-transition front]
\label{rem:agn-small-transition}
AGN's Section~II identifies three routes to the i.i.d.\ boundary, the first
being \(p\approx0\) or \(q\approx0\).  Their Theorem~1 and Corollary~1 give
non-asymptotic bounds with explicit dependence on the lower margin
\(\delta_{\rm A}\leq\min(p,q)\), so the small-transition front is covered by
their published result
\cite[Section~II, Theorem~1 and Corollary~1]{AbrahamGassiatNaulet2023}.

The precise scope still requires their spectral-gap qualification.  They work
on
\[
 \Theta_{L_{\rm A}}
 =\Theta(\delta_{\rm A},\epsilon_{\rm A},\zeta_{\rm A})
 \cap\{1-|1-p-q|\geq L_{\rm A}\},
\]
treat \(L_{\rm A}\) as fixed, and state in Remark~1 that this permits one of
\(p,q\) to vanish, but not both.  For the joint symmetric path
\(p=q=u\downarrow0\), membership requires \(L_{\rm A}\leq2u\); hence no single
fixed-\(L_{\rm A}\) class contains the whole path.  The point
\((u,u)\) is covered for every fixed \(u>0\) after choosing
\(L_{\rm A}\leq2u\).  Their theorem statement suppresses the dependence of
\(C(K,L_{\rm A})\), but their proof of Lemma~1 permits the explicit tracking
carried out below.  AGN themselves state that accurate rates without a
spectral-gap lower bound are beyond their scope and that tracking the gap with
their method leaves an upper/lower mismatch by a factor
\(\delta_{\rm A}\).
The distinction is visible directly from
\[
 \operatorname{Cov}\{\mathbf 1_{\{S_0=2\}},\mathbf 1_{\{S_h=2\}}\}
 =\frac{pq}{(p+q)^2}(1-p-q)^h.
\]
If \(p\downarrow0\) and \(q=q_0>0\), this covariance tends to zero; if
\(p=q=u\downarrow0\), it tends to \(1/4\) for every fixed \(h\).  Thus the
one-sided front approaches a one-component i.i.d.\ law, whereas simultaneous
slow mixing does not.
\end{remark}

We now quantify that dependence.  Following
\cite[proof of Lemma~1]{AbrahamGassiatNaulet2023}, set
\begin{equation}
 m_L=\left\lceil\frac{\log4}{L}\right\rceil+2,
 \qquad
 \gamma_L=\frac1{2m_L}.
 \label{eq:agn-explicit-gap}
\end{equation}
Their length-three window chain has mixing time at most \(m_L\), and the
pseudo-spectral-gap comparison used in their proof gives
\(\gamma_{\rm ps}\geq\gamma_L\).  Solving their displayed Bernstein
inequality coordinatewise and taking a union bound over the \(K^3\) cells
gives the admissible explicit radius
\begin{equation}
 r_3(T,y;L)
 \leq C_K\left\{
 \sqrt{\frac{y}{T\gamma_L}}
 +\frac{y}{T\gamma_L}
 \right\},
 \qquad
 y=\log\frac{2K^3}{\alpha}.
 \label{eq:agn-explicit-moment-radius}
\end{equation}
This is the dependence hidden in their fixed-\(L\) notation; it is an
explicit admissible tracking of their proof, not a claim that their constants
are optimal.

Fix \(\epsilon_0,\zeta_0>0\), put
\[
 \delta_{\rm A}=u,\qquad
 \epsilon_{\rm A}=\epsilon_0\leq|1-2u|,\qquad
 \zeta_{\rm A}=\zeta_0,\qquad L_{\rm A}=2u,
\]
and assume \(u<\epsilon_0\zeta_0\).  The inverse factor in AGN's
Theorem~1 is
\[
 \frac{\delta_{\rm A}\vee
 \epsilon_{\rm A}\zeta_{\rm A}}
 {\delta_{\rm A}\epsilon_{\rm A}^2\zeta_{\rm A}^3}
 =\frac1{u\epsilon_0\zeta_0^2}.
\]
Moreover,
\begin{equation}
 \gamma_{2u}
 =\frac1{2\{\lceil(\log4)/(2u)\rceil+2\}}
 \asymp u.
 \label{eq:agn-gap-symmetric}
\end{equation}
Substitution of \eqref{eq:agn-explicit-moment-radius} therefore yields the
fully tracked AGN sufficient radius
\begin{equation}
 r_Q^{\rm AGN}(u,T,\alpha)
 \leq C_{K,\epsilon_0,\zeta_0}
 \left\{
 \frac{\sqrt y}{u^{3/2}\sqrt T}
 +\frac{y}{u^2T}
 \right\}.
 \label{eq:agn-u-raw-rate}
\end{equation}
Thus tracking \(L_{\rm A}=2u\) adds the factor \(u^{-1/2}\) which was hidden
when \(C(K,L_{\rm A})\) was held fixed.

This is not yet a rate for \(v_\infty\).  At \(p=q=u\),
\begin{equation}
 v_\infty(u,a)=\sigma^2+a^2\frac{1-u}{u},\qquad
 \partial_pv_\infty=\partial_qv_\infty=-\frac{a^2}{2u^2},
 \qquad
 \frac{d}{du}v_\infty(u,a)=-\frac{a^2}{u^2}.
 \label{eq:agn-u-gradient}
\end{equation}
Fixing \(a\) and applying the ordinary delta upper bound in the local regime
\(r_Q^{\rm AGN}=o(u)\) therefore gives
\begin{align}
 r_{v,\rm delta}^{\rm AGN}
 &\leq
 \left(
 |\partial_pv_\infty|+|\partial_qv_\infty|
 \right)r_Q^{\rm AGN}
 +O\{(r_Q^{\rm AGN})^2/u^3\}
 \nonumber\\
 &=O\left\{
 a^2\left(
 \frac{\sqrt y}{u^{7/2}\sqrt T}
 +\frac{y}{u^4T}
 \right)
 \right\}.
 \label{eq:agn-u-delta-rate}
\end{align}
For \(Nu\gg1\), the same calculation gives
\(r_{G_N,\rm delta}^{\rm AGN}
=O(Nr_{v,\rm delta}^{\rm AGN})\).
The delta propagation adds another \(u^{-2}\).  The leading tracked AGN
functional exponent is therefore \(u^{-7/2}\), whereas the local long-lag
upper bound in Theorem~\ref{thm:persistence-local-upper} has exponent
\(u^{-3/2}\).

\begin{lemma}[Exact KL control on the symmetric boundary path]
\label{lem:agn-symmetric-exact-kl}
Let \(\overline P_u^{(T)}\) and \(P_u^{(T)}\) denote respectively the
complete-data and observed-data laws in the symmetric Markov model
\[
 Q_u=\begin{pmatrix}1-u&u\\u&1-u\end{pmatrix},
\]
with stationary initialization and the same fixed Gaussian emissions under
both alternatives.  If \(T\geq2\), \(0<u\leq1/4\), and \(0<d\leq u/2\), then
\begin{align}
 D_{\rm KL}\!\left(P_u^{(T)}\middle\|P_{u+d}^{(T)}\right)
 &\leq
 D_{\rm KL}\!\left(\overline P_u^{(T)}
 \middle\|\overline P_{u+d}^{(T)}\right)
 \nonumber\\
 &=(T-1)\operatorname{kl}(u\|u+d)
 \nonumber\\
 &\leq
 \frac{(T-1)d^2}{(u+d)(1-u-d)}
 \leq\frac{8(T-1)d^2}{5u}.
 \label{eq:agn-exact-complete-kl}
\end{align}
\end{lemma}

\begin{proof}
Put \(v=u+d\).  The stationary distribution is
\(\pi_u=\pi_v=(1/2,1/2)\), and the conditional Gaussian laws are unchanged.
The chain rule therefore gives the exact calculation
\begin{align*}
 D_{\rm KL}\!\left(\overline P_u^{(T)}
 \middle\|\overline P_v^{(T)}\right)
 &=D_{\rm KL}(\pi_u\|\pi_v)
 +\sum_{t=1}^{T-1}
 \mathbb E_u\!\left[
 \log\frac{Q_u(S_t,S_{t+1})}{Q_v(S_t,S_{t+1})}
 \right]
 \nonumber\\
 &\quad+
 \sum_{t=1}^{T}
 \mathbb E_u\!\left[
 \log\frac{g_{S_t}(X_t)}{g_{S_t}(X_t)}
 \right]
 \\
 &=(T-1)\left\{
 u\log\frac{u}{v}
 +(1-u)\log\frac{1-u}{1-v}
 \right\}
 \\
 &=(T-1)\operatorname{kl}(u\|v).
\end{align*}
Indeed, under either transition row the switch indicator is
\(\operatorname{Bern}(u)\).  Projection from the complete path onto the
observations and the data-processing inequality give the first inequality in
\eqref{eq:agn-exact-complete-kl}.  Finally,
\[
 \operatorname{kl}(u\|u+d)
 \leq
 \chi^2\{\operatorname{Bern}(u)\|
 \operatorname{Bern}(u+d)\}
 =\frac{d^2}{(u+d)(1-u-d)}.
\]
Since \(u+d\geq u\) and
\(1-u-d\geq5/8\), the final inequality follows.
\end{proof}

The Fisher expansion is an interpretation of
\eqref{eq:agn-exact-complete-kl}, not its justification.  If
\(f_u(d)=\operatorname{kl}(u\|u+d)\), then
\[
 f_u(0)=f_u'(0)=0,\qquad
 f_u''(0)=\frac1{u(1-u)},
\]
and Taylor's theorem gives, for \(0\leq d\leq u/2\),
\begin{equation}
 \operatorname{kl}(u\|u+d)
 =\frac{d^2}{2u(1-u)}
 +O\left(\frac{d^3}{u^2}\right).
 \label{eq:agn-kl-taylor-domain}
\end{equation}
More explicitly,
\[
 f_u'''(d)
 =-\frac{2u}{(u+d)^3}
 +\frac{2(1-u)}{(1-u-d)^3},
 \qquad
 \sup_{0\leq d\leq u/2}|f_u'''(d)|\leq \frac{C}{u^2},
\]
so the ratio of the Taylor remainder to the quadratic term is
\(O(d/u)\).  The quadratic approximation is therefore relatively accurate
only when \(d/u\to0\).  At the Fisher scale
\(d\asymp\sqrt{u(1-u)/T}\), this condition is \(Tu\to\infty\).
The earlier working draft stated the resulting
\(u^{-3/2}T^{-1/2}\) expression without this regime condition; that was an
invalid extrapolation of the Taylor approximation and is corrected below.

The lower bound itself can be made uniform without Taylor expansion.  Fix
\(0<c_0\leq1/2\) and set
\begin{equation}
 d_{T,u}
 =c_0\min\left\{
 u,\sqrt{\frac{u(1-u)}{T-1}}
 \right\}.
 \label{eq:agn-boundary-perturbation}
\end{equation}
Lemma~\ref{lem:agn-symmetric-exact-kl} gives
\[
 D_{\rm KL}\!\left(P_u^{(T)}\middle\|
 P_{u+d_{T,u}}^{(T)}\right)
 \leq\frac85c_0^2
 \min\{(T-1)u,1-u\}
 \leq\frac85c_0^2.
\]
Moreover, the functional separation is exact:
\[
 |v_\infty(u+d_{T,u},a)-v_\infty(u,a)|
 =\frac{a^2d_{T,u}}{u(u+d_{T,u})}.
\]
The two-point inequality therefore implies, for every estimator,
\begin{align}
 &\max_{\vartheta\in\{u,u+d_{T,u}\}}
 P_\vartheta^{(T)}
 \left\{
 |\widehat v_{\infty,T}-v_\infty(\vartheta,a)|
 \geq
 \frac{a^2d_{T,u}}{2u(u+d_{T,u})}
 \right\}
 \nonumber\\
 &\hspace{35mm}\geq
 \frac14\exp\left(-\frac85c_0^2\right).
 \label{eq:agn-boundary-two-regime-lb}
\end{align}
Since \(u+d_{T,u}\leq3u/2\), this construction yields the two-point
lower-bound scale
\begin{equation}
 a^2\min\left\{
 \frac1u,\,
 \frac{\sqrt{1-u}}{u^{3/2}\sqrt{T-1}}
 \right\}.
 \label{eq:agn-persistence-gap}
\end{equation}
Thus \(a^2u^{-3/2}T^{-1/2}\) is the \(Tu\gg1\) branch.  When
\(Tu=O(1)\), the legitimate boundary branch is \(a^2/u\), not an
extrapolation of the Fisher formula.

\begin{corollary}[Confidence-calibrated persistence lower radius]
\label{cor:agn-boundary-confidence}
Let \(0<\delta<1/4\),
\(\ell_\delta=\log\{1/(4\delta)\}\), and
\begin{equation}
 d_{T,u,\delta}
 =\min\left\{
 \frac u2,
 \sqrt{\frac{5u\ell_\delta}{8(T-1)}}
 \right\}.
 \label{eq:agn-confidence-perturbation}
\end{equation}
Then every estimator satisfies
\begin{align}
 &\max_{\vartheta\in\{u,u+d_{T,u,\delta}\}}
 P_\vartheta^{(T)}
 \left\{
 |\widehat v_{\infty,T}-v_\infty(\vartheta,a)|
 \geq
 \frac{a^2d_{T,u,\delta}}
 {2u(u+d_{T,u,\delta})}
 \right\}
 \geq\delta,
 \label{eq:agn-confidence-lower}
\end{align}
and the threshold on the left is at least
\begin{equation}
 \frac{a^2}{6}
 \min\left\{
 \frac1u,
 \frac{\sqrt{\ell_\delta}}
 {u^{3/2}\sqrt{T-1}}
 \right\}.
 \label{eq:agn-confidence-lower-scale}
\end{equation}
\end{corollary}

\begin{proof}
Equation~\eqref{eq:agn-exact-complete-kl} and
\eqref{eq:agn-confidence-perturbation} give
\[
 D_{\rm KL}\!\left(P_u^{(T)}\middle\|
 P_{u+d_{T,u,\delta}}^{(T)}\right)
 \leq\ell_\delta.
\]
Bretagnolle--Huber therefore gives
\((1/4)e^{-\ell_\delta}=\delta\).  Since
\(u+d_{T,u,\delta}\leq3u/2\), the functional threshold is at least
\(a^2d_{T,u,\delta}/(3u^2)\), which yields
\eqref{eq:agn-confidence-lower-scale}.
\end{proof}

For \(u\leq1/6\), take \(u_\star=u\), \(\alpha=1\), and \(\beta=3/2\) in
Theorem~\ref{thm:persistence-local-upper}; the lower pair in
Corollary~\ref{cor:agn-boundary-confidence} lies in that same local cell.
For fixed known \(a,\sigma\) with non-degenerate signal-to-noise ratio, the
upper and lower radii have identical powers:
\[
 \begin{array}{c|c|c}
 &\text{lower radius}&\text{upper radius}\\\hline
 Tu\gg\log(1/\delta)
 &a^2u^{-3/2}\sqrt{\log(1/\delta)/T}
 &( |a|+\sigma)^2u^{-3/2}
   \sqrt{\log(1/\delta)/T}\\
 Tu=O\{\log(1/\delta)\}
 &a^2u^{-1}&a^2u^{-1}
 \end{array}
\]
Thus the persistent-path result is minimax-rate tight on a local cell known
up to a fixed multiplicative factor.  The regular branch is a genuine
statistical matching result.  The boundary branch is diameter saturation and
does not imply consistent relative estimation.  This is a local minimax
statement indexed by the prescribed scale \(u_\star\).  Since the estimator
itself depends on \(u_\star\), it does not imply a single estimator attaining
these radii uniformly over unknown persistence scales.  No global adaptive
minimax claim is made.

This calculation does not assume local asymptotic normality of the observed
likelihood.  It uses an exact complete-data identity followed by data
processing, so any singularity of the observed HMM can only make the observed
KL smaller.  The price is possible looseness: a complete-data KL upper bound
need not reveal the sharp observed-data geometry.  AGN's reparametrization and
Proposition~2 are needed precisely to control that sharper marginal geometry
near an i.i.d.\ boundary.  In their notation
\(\phi_2=1-p-q\); Proposition~2 assumes \(|\phi_2|\) small, whereas on the
symmetric slow-mixing path \(\phi_2=1-2u\to1\).  The two KL calculations
therefore concern different degeneracies.

Finally, \(d_{T,u}\downarrow0\) is a separation used to prove an impossibility
result, not by itself an achievable estimation radius.  In the \(Tu\gg1\)
branch its relative size is
\[
 \frac{d_{T,u}}u
 \asymp\frac1{\sqrt{Tu}},
\]
which worsens as \(u\) decreases.  Omitting confidence and fixed
signal-to-noise factors, the explicit comparison is
\[
 \begin{array}{c|c|c}
 &\text{radius for }u&\text{radius for }v_\infty\\\hline
 \text{AGN tracked uniform upper}
 &u^{-3/2}T^{-1/2}&a^2u^{-7/2}T^{-1/2}\\
 \text{long-lag local upper, }Tu\gg1
 &u^{1/2}T^{-1/2}&a^2u^{-3/2}T^{-1/2}\\
 \text{two-point local lower, }Tu\gg1
 &u^{1/2}T^{-1/2}&a^2u^{-3/2}T^{-1/2}\\
 \text{local boundary, }Tu=O(1)
 &u&a^2u^{-1}
 \end{array}
\]
The tracked AGN radius is looser by \(u^{-2}\) on the symmetric Gaussian
submodel.  Their proof first estimates a fixed length-three law at the
small-gap scale \((Tu)^{-1/2}\) and then applies an inverse bound of order
\(u^{-1}\).  In contrast, the long-lag moment has the same
\((Tu)^{-1/2}\) fluctuation but its inverse derivative is of order \(u\),
giving \(\sqrt{u/T}\).  The discrepancy is therefore attributable to the
fixed-block uniform AGN procedure and its broader parameter class, not to an
inherent lower bound on the present local Gaussian subexperiment.

Nor can \eqref{eq:agn-u-raw-rate} be extrapolated at fixed \(T\) all the way
to \(u=0\): AGN's stated confidence range requires
\(y\leq T u^2\epsilon_0^4\zeta_0^6\), and their theorem treats
\(L_{\rm A}\) as fixed, while the symmetric path has \(L_{\rm A}=2u\).
Theorem~\ref{thm:persistence-local-upper} supplies the matching observed-data
upper bound for \eqref{eq:agn-persistence-gap} on a prescribed local
multiplicative cell.  This comparison supplies no global adaptive theorem
because the lag scale is fixed in advance.

For comparison, the fixed-\(L_{\rm A}\) small-transition front covered
directly by AGN may be represented by \(p=u\downarrow0\) and
\(q=q_0>0\) fixed.  Then \(L_{\rm A}\) remains bounded away from zero and
\[
 v_\infty(p,q_0,a)-\sigma^2
 =\frac{4a^2p q_0(2-p-q_0)}{(p+q_0)^3}
 =\frac{4a^2(2-q_0)}{q_0^2}p+O(p^2).
\]
Here the functional derivative stays finite, so the AGN
\(\delta_{\rm A}^{-1}T^{-1/2}\) raw radius propagates with the same
\(\delta_{\rm A}^{-1}\) power.  This is distinct from sending both hazards
to zero.

Finally, for finite \(D\), Markovization produces \(2D\) latent age states
with identical emission laws across the \(D\) ages of each regime.  AGN's
two-state theorem does not apply to that augmented representation, and the
repeated emission columns preclude treating it as an unrestricted identifiable
\(2D\)-state HMM.

\paragraph{Discounted Markov-cost variance.}
To avoid the collision of notation with this article, write \(B\) for the
total transition budget, \(H\) for the truncation horizon, and
\(\rho\in[0,1)\) for the discount factor in
\cite{ThoppePrashanthBhat2024}.  Their target and truncated return are
\[
 D_\rho=\sum_{t=0}^\infty\rho^tf(S_t),\qquad
 D_{\rho,H}=\sum_{t=0}^{H-1}\rho^tf(S_t).
\]
Their reset access produces
\(m=\lceil B/H\rceil\) independent copies of \(D_{\rho,H}\); states and
costs are observed, and \(|f(s)|\leq K\).  In the published version,
Theorem~4.12 states
\begin{align}
 \mathbb E|\widehat V_B-\operatorname{Var}(D_\rho)|
 &\leq
 \frac{8(1-\rho^H)^2K^2}{\sqrt m(1-\rho)^2}
 +\frac{4\rho^HK^2}{(1-\rho)^2},
 \label{eq:tpb-expectation}\\
 \mathbb P\{|\widehat V_B-\operatorname{Var}(D_\rho)|>\varepsilon\}
 &\leq
 2\exp\left[
 -\frac{m(1-\rho)^4}{32(1-\rho^H)^4K^4}
 \left\{\varepsilon-
 \frac{4\rho^HK^2}{(1-\rho)^2}\right\}^2
 \right]
 \label{eq:tpb-tail}
\end{align}
for
\(\varepsilon>4\rho^HK^2/(1-\rho)^2\).  Equivalently, its
\((1-\alpha)\)-confidence radius is
\begin{equation}
 r_{\rm TPB}(B,H,\alpha)
 =\frac{4\rho^HK^2}{(1-\rho)^2}
 +\frac{\sqrt{32}(1-\rho^H)^2K^2}{(1-\rho)^2}
 \sqrt{\frac{\log(2/\alpha)}m}.
 \label{eq:tpb-radius}
\end{equation}
Thus Theorem~4.12 is an upper concentration result for an episodic sample
variance, not a lower bound for a hidden single-trajectory experiment.

The comparable horizon limit in our notation is \(N\to\infty\), not
\(T\to\infty\): our \(T\) is data length, whereas their \(H\), denoted \(T\)
in their paper, is an estimator truncation parameter.  If
\(c_h=\operatorname{Cov}(X_0,X_h)\) and
\(\sum_{h\geq0}|c_h|<\infty\), then
\begin{align}
 V_\rho
 &:={\operatorname{Var}}\left(\sum_{t=0}^\infty\rho^tX_t\right)
 =\frac{c_0+2\sum_{h=1}^\infty\rho^hc_h}{1-\rho^2},
 \label{eq:tpb-abelian-exact}\\
 (1-\rho^2)V_\rho&\longrightarrow v_\infty,
 &
 \frac{G_N}{N}&\longrightarrow v_\infty.
 \label{eq:tpb-abelian-limit}
\end{align}
Hence, with the effective horizon
\(N_\rho=(1-\rho^2)^{-1}\), one has the population equivalence
\(V_\rho\sim G_{N_\rho}\).  This is the only direct convergence between
the two formulations.

Indeed, set \(h=1-\rho\) and normalize the discounted target as
\(W_\rho=(1-\rho^2)V_\rho\).  Multiplying
\eqref{eq:tpb-radius} by \(1-\rho^2\) gives
\begin{align}
 r_{W}(B,H,\alpha)
 &=\frac{4\rho^HK^2(1+\rho)}h
 \nonumber\\
 &\quad+
 \frac{\sqrt{32}(1-\rho^H)^2K^2(1+\rho)}h
 \sqrt{\frac{\log(2/\alpha)}m}.
 \label{eq:tpb-normalized-radius}
\end{align}
To make the first term at most \(\eta/2\), it is necessary that
\begin{equation}
 \rho^H\leq\frac{\eta h}{8K^2(1+\rho)},\qquad
 H\geq
 \frac{\log\{8K^2(1+\rho)/(\eta h)\}}{-\log\rho}
 \sim\frac{\log(1/h)}h.
 \label{eq:tpb-truncation-need}
\end{equation}
Once this truncation bias is controlled, the stochastic term in
\eqref{eq:tpb-normalized-radius} requires
\[
 m\gtrsim
 \frac{K^4}{\eta^2h^2}\log\frac2\alpha.
\]
Since \(B\asymp mH\), the resulting sufficient budget obeys
\begin{equation}
 B\gtrsim
 \frac{K^4}{\eta^2h^3}
 \log\frac2\alpha
 \log\frac{K^2}{\eta h}.
 \label{eq:tpb-singular-budget}
\end{equation}
Thus the normalized Theorem~4.12 bound is singular as
\(\rho\uparrow1\) and does not converge to the
\(\sqrt{\log(1/\alpha)/T}\) bound for \(v_\infty\).  The discrepancy
survives the population normalization because TPB uses independent reset
episodes, a generic bounded-cost sample-variance estimator, and an explicit
truncation bias, whereas the present experiment consists of one stationary
hidden trajectory with unbounded Gaussian emissions and an exact parametric
functional.  Their general lower bound already establishes the root-sample
barrier for variance risk; consequently that power is not new here either.
The present result is therefore not merely a finite-horizon extension of
Theorem~4.12.  Its distinct scope is the explicit functional separation for
\(G_N\) and \(v_\infty\) in the stationary hidden capped-duration experiment.
\subsection{Scope of the optimality claim}
\label{subsec:optimality-scope}

On the fixed interior classes, the upper and lower bounds match at the
minimax-rate level; the global Doeblin and inversion constants are not claimed
sharp.  Along the persistence path,
Theorem~\ref{thm:persistence-local-upper} and
Corollary~\ref{cor:agn-boundary-confidence} match powers cell by cell, but the
lag uses the prescribed scale $u_\star$.  We do not prove adaptation over a
union of unknown scales.  For unbounded durations controlled only through two
moments, we make no uniform exponential claim; such a result would require
stronger tail assumptions than are imposed here.
